\documentclass[12pt]{article}
\usepackage[utf8]{inputenc}
\usepackage[T1]{fontenc}
\usepackage[margin=1in]{geometry}
\usepackage{setspace}
\usepackage{xcolor}
\usepackage[colorlinks=true, urlcolor=blue!60!black, linkcolor=blue!60!black]{hyperref}
\setlength{\parindent}{2em}
\setlength{\parskip}{1.25em}
\renewcommand{\baselinestretch}{1.0}
\begin{document}
\begin{center}
{\LARGE \textit{Barron v.\ Mayor \& City Council of Baltimore}, 32 U.S.\ (7 Pet.) 243 (1833)\par}
\vspace{0.8em}
{\large Case Brief --- Four-Page Edition\par}
\end{center}
\vspace{1em}

\section{Memory Jogger}

Chief Justice John Marshall held for a unanimous Court that the Fifth Amendment's Takings Clause---and by extension the entire federal Bill of Rights---restrains only the federal government and does not apply to state or municipal governments (\href{https://supreme.justia.com/cases/federal/us/32/243/}{Justia}).

\section{Detailed Case Facts}

John Barron, surviving partner of John Craig, owned a profitable commercial wharf in eastern Baltimore that enjoyed the deepest water in the harbor. Between 1815 and 1821, Baltimore regraded and paved streets under its charter and Maryland statutes, diverting several natural streams so that they emptied toward the wharf. Heavy rains washed sand and earth down from the hills and deposited them in front of the wharf, making the water so shallow that large vessels could no longer reach it and rendering it nearly worthless. The city paid no compensation. Barron contended that the city's state-authorized destruction of his wharfage rights violated the Fifth Amendment's command that ``private property [shall not] be taken for public use, without just compensation.'' The Fourteenth Amendment did not yet exist.

\section{Procedural History}

In 1822 Barron sued the city in the Baltimore County Court, which ruled against the city on all points; a jury awarded him \$4,500. The Court of Appeals for the Western Shore of Maryland reversed. Barron took a writ of error to the U.S. Supreme Court under Section 25 of the Judiciary Act of 1789, which allowed review of state-court decisions upholding state authority against a federal constitutional claim. The Court heard argument in February 1833 and dismissed the writ for want of jurisdiction on February 16, 1833.

\section{Judicial Votes}

Writ dismissed, \textbf{7--0}.
\begin{itemize}
\item \textbf{Majority (7):} Chief Justice Marshall, joined by Justices Johnson, Duvall, Story, Thompson, McLean, and Baldwin.
\item \textbf{Concurrences and dissents:} None.
\end{itemize}

\section{Holding}

The Takings Clause of the Fifth Amendment---and, by the opinion's reasoning, the Bill of Rights generally---limits only the federal government and does not bind state or local governments. Because no federal constitutional question was presented, the Court lacked jurisdiction. Strictly, the holding was confined to the Just Compensation Clause; Marshall's broader statement that ``[t]hese amendments contain no expression indicating an intention to apply them to the State governments'' functioned as sweeping dicta that later antebellum cases treated as settled. \textit{Barron} has never been formally overruled, but its practical effect was largely superseded by selective incorporation through the Fourteenth Amendment, beginning with the Takings Clause in \textit{Chicago, Burlington \& Quincy Railroad Co. v. Chicago} (1897).

\section{Analysis of Opinions}

\textbf{Opinion of the Court (Marshall).} Marshall called the issue ``of great importance, but not of much difficulty'' and gave four linked arguments. \emph{Structure:} the Constitution was ordained by the people ``for their own government, and not for the government of the individual States,'' and each State had its own constitution with its own limits. \emph{Intratextual comparison:} Article I, \S~9 uses general language that binds only Congress, while \S~10 repeats the same prohibitions with express ``No State shall'' language; whenever the Framers meant to bind the States they said so ``in positive words,'' and the Fifth Amendment does not. \emph{Article V logic:} citizens wanting protection from their own legislature could amend their state constitution; ``[t]he unwieldy and cumbrous machinery'' of federal amendment would never have been chosen for that purpose. \emph{Ratification history:} the Bill of Rights answered Anti-Federalist fears of the new national government, ``not against those of the local governments.''

\textbf{Separate opinions.} None.

\section{Examples: Future Applications}

\begin{enumerate}
\item \textbf{Same-Side: State Prosecution by Information.} A State charges murder by a prosecutor's information instead of a grand jury indictment. Permitted: the Grand Jury Clause does not bind the States under \textit{Barron}, and \textit{Hurtado v. California} (1884) declined to incorporate it.
\item \textbf{Same-Side: No Civil Jury in State Court.} A State abolishes civil jury trials in tort suits. The Seventh Amendment has never been incorporated, so the State may structure civil trials under its own constitution.
\item \textbf{Opposite-Side: Taking for \$1.} A city condemns a railroad right-of-way for a street and awards \$1. \textit{Barron} bars direct reliance on the Fifth Amendment, but \textit{Chicago, Burlington \& Quincy} held such a taking violates Fourteenth Amendment due process.
\item \textbf{Opposite-Side: Handgun Ban or Excessive Forfeiture.} A city bans handguns in the home, or a State forfeits a \$42,000 vehicle for a minor drug offense. \textit{McDonald} (2010) and \textit{Timbs} (2019) incorporated the Second Amendment and the Excessive Fines Clause.
\item \textbf{Fence-Sitter: State Troops Quartered in Homes.} A governor orders non-federalized Guard troops into private homes. The Supreme Court has never decided whether the Third Amendment is incorporated (the Second Circuit held it was in \textit{Engblom v. Carey}); a court would ask whether the right is ``deeply rooted in this Nation's history and tradition.''
\end{enumerate}

\section{Critique}

\textbf{Originalist and textualist views.} Most originalists praise \textit{Barron}'s disciplined intratextualism, corroborated by the First Congress's rejection of Madison's proposed state-binding amendment. William Winslow Crosskey countered that only the First Amendment names ``Congress,'' implying that the others bound the States too, a thesis Charles Fairman disputed. Akhil Amar and Kurt Lash stress that John Bingham cited \textit{Barron} in 1866 to explain why the Fourteenth Amendment needed ``No State shall'' language; on this view the real error was the \textit{Slaughter-House Cases}' gutting of the Privileges or Immunities Clause.

\textbf{Progressive views and logical weaknesses.} Marshall's assurance that citizens could simply reform their state constitutions ignored that enslaved people, free Black citizens, women, and dissenters lacked political power; \textit{Barron} left Southern States free to suppress abolitionist speech until incorporation began with \textit{Gitlow} (1925). Marshall also glossed over \S~10's Contract, Bill of Attainder, and Ex Post Facto Clauses, which protect individual rights, and treated the Bill of Rights as a bargain between sovereigns rather than a declaration of pre-existing rights.

\section{Key Quotations}

\begin{quote}
``The powers they conferred on this government were to be exercised by itself, and the limitations on power, if expressed in general terms, are naturally, and we think necessarily, applicable to the government created by the instrument. They are limitations of power granted in the instrument itself, not of distinct governments framed by different persons and for different purposes.''
\end{quote}
\noindent --- Chief Justice Marshall, Opinion of the Court, 32 U.S. at 247.

\begin{quote}
``These amendments demanded security against the apprehended encroachments of the General Government --- not against those of the local governments. In compliance with a sentiment thus generally expressed, to quiet fears thus extensively entertained, amendments were proposed by the required majority in Congress and adopted by the States. These amendments contain no expression indicating an intention to apply them to the State governments. This court cannot so apply them.''
\end{quote}
\noindent --- Chief Justice Marshall, Opinion of the Court, 32 U.S. at 250.

\begin{quote}
``The question thus presented is, we think, of great importance, but not of much difficulty. The Constitution was ordained and established by the people of the United States for themselves, for their own government, and not for the government of the individual States.''
\end{quote}
\noindent --- Chief Justice Marshall, Opinion of the Court, 32 U.S. at 247.

\end{document}
